\chapter{Petle: powtarzanie pracy bez kopiowania kodu}

\section{Po co istnieja petle}
Petla pozwala wykonac podobna operacje wiele razy. To jedna z najwazniejszych idei
programowania. Zamiast recznie pisac dziesiec podobnych sekcji HTML albo dziesiec razy
zwiekszac licznik, mozna zapisac regule powtarzania.

Na stronach internetowych petle przydaja sie np. do:

\begin{itemize}
  \item generowania list wpisow,
  \item budowania siatki kart,
  \item przechodzenia po elementach menu,
  \item wykonywania operacji na wszystkich rekordach danych,
  \item tworzenia wielu podstron wedlug jednego schematu.
\end{itemize}

MoonChunk wspiera dwa glowne rodzaje petli: \mc{for} i \mc{while}.

\section{Petla \mc{for}}
Petla \mc{for} jest bardzo dobra wtedy, gdy chcesz wyraznie kontrolowac licznik.
W MoonChunk ma ona klasyczny, jawny zapis.

\begin{lstlisting}[style=moonchunk,caption={Klasyczna petla for}]
for (let i: int = 0; i < 5; i++) {
  print(i);
};
\end{lstlisting}

Ten zapis czyta sie tak:

\begin{enumerate}
  \item utworz licznik \mc{i} i ustaw go na \mc{0},
  \item dopoki \mc{i < 5}, wykonuj ciało petli,
  \item po kazdym obiegu zwieksz \mc{i} o jeden.
\end{enumerate}

\section{Budowa petli \mc{for} krok po kroku}
Petla \mc{for} ma trzy glowne czesci wewnatrz nawiasow:

\begin{itemize}
  \item inicjalizacje,
  \item warunek,
  \item aktualizacje.
\end{itemize}

\begin{lstlisting}[style=moonchunk,caption={Trzy czesci petli for}]
for (let i: int = 0; i < 10; i = i + 2) {
  print(i);
};
\end{lstlisting}

Tutaj licznik rosnie co dwa kroki, a nie co jeden.

\section{Czy aktualizacja w \mc{for} musi byc \mc{i++}}
Nie. MoonChunk pozwala wpisac zwykle wyrazenie aktualizujace. Dlatego poprawne sa oba
style:

\begin{lstlisting}[style=moonchunk,caption={Dwie formy aktualizacji licznika}]
for (let i: int = 0; i < 10; i++) {
  print(i);
};

for (let i: int = 0; i < 10; i = i + 2) {
  print(i);
};
\end{lstlisting}

To bardzo praktyczne, bo pozwala zapisac nie tylko krok o jeden, ale dowolna logike
przesuwania licznika.

\section{Petla \mc{while}}
\mc{while} przydaje sie wtedy, gdy nie chcesz mysliec wprost o klasycznym liczniku
albo liczba powtorzen nie jest z gory oczywista. Skladnia jest prosta:

\begin{lstlisting}[style=moonchunk,caption={Podstawowa petla while}]
let counter: int = 0;

while (counter < 3) {
  counter = counter + 1;
};
\end{lstlisting}

Petla powtarza swoje ciało tak dlugo, jak warunek pozostaje \mc{true}.

\section{Warunek petli musi byc typu \mc{bool}}
Podobnie jak przy \mc{if}, MoonChunk wymaga, by warunek w \mc{for} i \mc{while}
byl wartoscia logiczna. To oznacza, ze zapis taki jak:

\begin{lstlisting}[style=moonchunk,caption={Niepoprawny warunek petli}]
while (3) {
  print("blad");
};
\end{lstlisting}

powinien skonczyc sie bledem typu.

\section{Petle a generowanie HTML}
Petle sa szczegolnie ciekawe wtedy, gdy budujesz wiele podobnych fragmentow strony.
MoonChunk pozwala umieszczac \mc{content} wewnatrz petli, jesli znajdujesz sie w ciele
\mc{page}.

\begin{lstlisting}[style=moonchunk,caption={Content generowany w petli}]
page "" {
  const cards = [
    { title: "A", description: "one", button: "Read more" },
    { title: "B", description: "two", button: "Read more" },
    { title: "C", description: "three", button: "Read more" }
  ];

  for (let i: int = 0; i < 3; i++) {
    content {
      <article class="card">
        <h3>{cards[i].title}</h3>
        <p>{cards[i].description}</p>
        <button>{cards[i].button}</button>
      </article>
    };
  };
};
\end{lstlisting}

To jeden z najbardziej naturalnych wzorcow pracy z MoonChunk.

\section{\mc{break} i \mc{continue}}
W petlach dostepne sa dwie bardzo wazne instrukcje sterujace.

\subsection*{\mc{break}}
Natychmiast konczy cala petle.

\subsection*{\mc{continue}}
Pomija reszte biezacego obiegu i przechodzi do kolejnej iteracji.

\begin{lstlisting}[style=moonchunk,caption={break i continue w praktyce}]
for (let i: int = 0; i < 10; i++) {
  if (i % 2 != 0) {
    continue;
  };

  if (i == 6) {
    break;
  };

  print(i);
};
\end{lstlisting}

W tym przykładzie:

\begin{itemize}
  \item liczby nieparzyste sa pomijane,
  \item cala petla konczy sie przy \mc{i == 6}.
\end{itemize}

\section{Gdzie wolno uzywac \mc{break} i \mc{continue}}
Tylko wewnatrz petli. To bardzo wazna zasada. Jesli sprobujesz uzyc tych instrukcji
poza \mc{for} albo \mc{while}, MoonChunk zglosi blad sterowania przeplywem.

\section{Petle zagniezdzone}
MoonChunk pozwala umieszczac jedna petle w drugiej.

\begin{lstlisting}[style=moonchunk,caption={Petle zagniezdzone}]
for (let i: int = 1; i <= 3; i++) {
  for (let j: int = 1; j <= i; j++) {
    print(i, j);
  };
};
\end{lstlisting}

To przydatne przy generowaniu siatek, tabel i struktur dwuwymiarowych.

\section{Petle a zakres nazw}
Ciało petli w MoonChunk korzysta z dosc rygorystycznych zasad redeklaracji. W praktyce
oznacza to, ze nie powinienes traktowac \mc{for} ani \mc{while} jak wolnej przestrzeni
na ponowne deklarowanie nazw z najblizszego otoczenia.

Jesli potrzebujesz celowo stworzyc osobny zakres do eksperymentalnej pracy z nazwa,
lepiej uzyc jawnego bloku \mc{\{ ... \}} albo wydzielic funkcje.

\section{Najczestsze bledy przy petlach}
\begin{itemize}
  \item zapomnienie o zmianie wartosci sterujacej w \mc{while},
  \item warunek niebedacy typu \mc{bool},
  \item uzycie \mc{break} albo \mc{continue} poza petla,
  \item bledna aktualizacja licznika,
  \item redeklaracja tej samej nazwy w nieodpowiednim zakresie,
  \item nieuwazne dzielenie albo modulo przez zero przy logice sterujacej petla.
\end{itemize}

\section{Kiedy wybrac \mc{for}, a kiedy \mc{while}}
Najprostsza zasada:

\begin{itemize}
  \item wybierz \mc{for}, gdy masz licznik lub jasny schemat iteracji,
  \item wybierz \mc{while}, gdy petla ma dzialac dopoki zachodzi pewien stan.
\end{itemize}

Jesli generujesz kolejne karty z tablicy wedlug indeksu, \mc{for} zwykle bedzie
bardziej naturalny. Jesli czekasz, az jakas wartosc osiagnie okreslony stan,
\mc{while} bywa prostszy.
